% Lec 07：有限生成阿贝尔群、同构定理续与群作用
\documentclass[UTF8,a4paper,11pt]{ctexart}
\usepackage[a4paper,margin=2.3cm]{geometry}
\usepackage{amsmath,amssymb,mathtools}
\usepackage{graphicx}
\usepackage{tikz}
\usetikzlibrary{arrows.meta}
\usepackage{enumitem}
\usepackage[hidelinks]{hyperref}
\usepackage{fancyhdr}
\setlist[itemize]{leftmargin=2.2em,itemsep=0.25em,topsep=0.35em}
\setlength{\parindent}{2em}
\setlength{\parskip}{0.35em}
\pagestyle{fancy}\fancyhf{}
\fancyhead[L]{离散数学与结构（I）}\fancyhead[R]{Lec 07}
\fancyfoot[C]{\thepage}\setlength{\headheight}{14pt}
\title{有限生成阿贝尔群、同构定理续与群作用}
\author{zhangjunshuo}\date{2026 年 9 月}
\begin{document}
\maketitle
\begin{quote}
本讲根据 2026 年 9 月 29 日的两份课堂转写、手写原稿和教材整理。转写中的口语识别错误已按手写稿与 Dummit--Foote《Abstract Algebra》（[DF]）和《离散数学与结构》（[DS]）校正。有限生成阿贝尔群部分保留课堂证明提纲，不把提纲写成完整的 Smith 标准形证明。
\end{quote}
\tableofcontents\clearpage

\section*{记号速查}
\begin{itemize}
\item $G,H,K$：群；$e$：单位元；$|G|$：群的阶；$H\le G$：子群；$N\trianglelefteq G$：正规子群。
\item $\langle S\rangle$：由 $S$ 生成的最小子群；$\langle g\rangle$：循环子群。
\item $G\times H$：外部直积；$G=H\times K$（内部直积）表示 $(h,k)\mapsto hk$ 是同构。
\item $\mathbb Z_n=\mathbb Z/n\mathbb Z$；$\mathbb Z_n^\times$：模 $n$ 的单位群；$\mathbb Z^r$：自由阿贝尔部分。
\item $Z(G)$：中心；$C_G(A)$：集合 $A$ 的中心化子；$c_g(x)=gxg^{-1}$：共轭映射。
\item $\operatorname{Hom}(G,H)$、$\operatorname{End}(G)$、$\operatorname{Aut}(G)$、$\operatorname{Inn}(G)$：同态、自同态、自同构、内自同构集合或群。
\item $G\curvearrowright X$：$G$ 作用在 $X$ 上；$G\cdot x$：轨道；$G_x$：稳定子；$[G:G_x]$：指数。
\item 二面体群按 [DF] 记 $D_{2n}$，阶为 $2n$；正方形的对称群写作 $D_8$。
\end{itemize}
\clearpage

\section{有限生成阿贝尔群与直积}


\textbf{本节用到的记号，先一次交代清楚}（第 6 讲已逐个定义过，这里只复述结论，方便单独阅读本节）：

\begin{itemize}
\item $G$：群；$e$：单位元；$|G|$：群的阶；$H\le G$：子群。
\item $S\subseteq G$：$G$ 的一个子集；$\langle S\rangle=\bigcap_{S\subseteq H\le G}H$：由 $S$ 生成的子群，即包含 $S$ 的所有子群的交。
\item $G=\langle g\rangle$：由单个元素生成，此时 $G$ 是\textbf{循环群}。
\item $C_n,\mathbb Z_n$：$n$ 阶循环群与整数模 $n$ 的剩余类群；$\mathbb Z$ 是无限循环群。
\item $\varphi:G\to H$：群同态；$\ker\varphi$：核；$\operatorname{im}\varphi$：像；$G/\ker\varphi\cong\operatorname{im}\varphi$。
\item $N\trianglelefteq G$：正规子群；$G/N$：商群；$\cong$：同构。
\item $G_1\times G_2$：群的直积，按分量运算 $(g_1,g_2)(h_1,h_2)=(g_1h_1,g_2h_2)$；$H,K\le G$ 且 $H\cap K=\{e\}$、$HK=G$ 时称\textbf{内部直积}。
\item $\gcd(m,n)$：最大公因数；$\mathbb Z_m\times\mathbb Z_n\cong\mathbb Z_{mn}\iff\gcd(m,n)=1$。
\item $|G|$ 的素幂分解：有限生成阿贝尔群结构定理写成 $\mathbb Z^r\times\mathbb Z_{p_1^{a_1}}\times\cdots$。
\end{itemize}


\subsection{生成子群与循环群}


设 $G$ 是群，$S\subseteq G$。由 $S$ 生成的子群记为 $\langle S\rangle$，它是包含 $S$ 的所有子群的交：

\[
\langle S\rangle=\bigcap_{S\subseteq H\le G}H.
\]

如果存在有限集 $S$ 使得 $G=\langle S\rangle$，就说 $G$ 是\textbf{有限生成群}。这和“有限群”不同：$\mathbb Z$ 是一个无限的有限生成群，而 $(\mathbb R,+)$ 不是有限生成群。

若 $G=\langle g\rangle$，则 $G$ 称为循环群。考虑

\[
\varphi_g:\mathbb Z\longrightarrow G,\qquad n\longmapsto g^n.
\]

这是一个满同态。若 $g$ 的阶为有限值 $m$，则

\[
\ker\varphi_g=m\mathbb Z,
\qquad G\cong\mathbb Z/m\mathbb Z=\mathbb Z_m.
\]

若 $g$ 的阶为无穷，则 $\ker\varphi_g=\{0\}$，从而 $G\cong\mathbb Z$。因此循环群恰好分为 $\mathbb Z$ 和 $\mathbb Z_m$ 两类；循环群都是阿贝尔群。

\subsection{直积与内部直积}


给定群 $G_1,G_2$，它们的外部直积为

\[
G_1\times G_2=\{(g_1,g_2):g_i\in G_i\},
\]

运算按分量定义：

\[
(g_1,g_2)(h_1,h_2)=(g_1h_1,g_2h_2).
\]

单位元是 $(e_1,e_2)$，逆元是 $(g_1^{-1},g_2^{-1})$。有限个群的直积递归地定义。

若 $H,K\le G$ 满足

\[
H\trianglelefteq G,\qquad K\trianglelefteq G,\qquad H\cap K=\{e\},\qquad HK=G,
\]

则称 $G$ 是 $H,K$ 的内部直积，并且

\[
G\cong H\times K,
\qquad (h,k)\longmapsto hk
\]

是一个同构。课堂证明中常用的等价条件是：$H\cap K=\{e\}$、$HK=G$，且 $hk=kh$ 对所有 $h\in H,k\in K$ 成立；存在性来自 $HK=G$，唯一性来自 $H\cap K=\{e\}$；交换性保证乘法映射保持群运算。

一个常用判据是：若 $\gcd(m,n)=1$，则

\[
\mathbb Z_m\times\mathbb Z_n\cong\mathbb Z_{mn};
\]

若 $\gcd(m,n)>1$，则直积中任一元素的阶都不能达到 $mn$，所以它不是循环群。

\subsection{有限生成阿贝尔群结构定理}


\textbf{结构定理（初等因子形式）}。每个有限生成阿贝尔群 $G$ 都同构于

\[
\mathbb Z^r\times
\mathbb Z_{p_1^{a_1}}\times\cdots\times
\mathbb Z_{p_t^{a_t}},
\]

其中 $r\ge 0$，每个 $p_i$ 是素数，$a_i\ge 1$。有限群时 $r=0$；自由秩 $r$ 以及素数幂因子的多重集合由 $G$ 唯一决定（因子次序除外）。

等价地，可以使用不变因子形式（这里采用 [DS] 常见的升序整除链记法；[DF] §5.2 常按降序编号，但内容相同）：

\[
G\cong \mathbb Z^r\times\mathbb Z_{d_1}\times\cdots\times\mathbb Z_{d_s},
\qquad d_1\mid d_2\mid\cdots\mid d_s,
\]

其中 $d_i>1$。这两种写法描述同一分类：前者把每个素数幂分开，后者把不同素数幂重新组合成整除链。

这里的“有限生成”是必要条件。比如 $(\mathbb R,+)$ 中任意有限生成子群都是可数的，而 $\mathbb R$ 不可数，因此 $\mathbb R$ 不可能由有限多个元素生成。

\subsection{课堂证明提纲：从一个最小关系开始}


课堂上给出的证明思路是对生成元个数归纳。设

\[
G=\langle a_1,\ldots,a_t\rangle
\]

是有限生成阿贝尔群，并考虑满同态

\[
\Phi:\mathbb Z^t\longrightarrow G,
\qquad (c_1,\ldots,c_t)\longmapsto c_1a_1+\cdots+c_ta_t.
\]

若 $\ker\Phi=\{0\}$，则 $G\cong\mathbb Z^t$。若核非零，则课堂上的“最小关系”有一个具体含义：在所有 $t$ 元生成元组 $(a_1,\ldots,a_t)$ 以及相应的非零整数关系 $(c_1,\ldots,c_t)\in\ker\Phi$ 中，按 $\sum_i|c_i|$ 取最小值；交换生成元、变号和做初等整数变换仍然属于同一类选择。音频中“把两个系数相减”指的正是这种生成元变换，和欧几里得算法的步骤相同。

通过可逆的整数初等变换（也就是把生成元换成另一组生成元 $b_1,\ldots,b_t$），可以把关系化为只有一个坐标非零：

\[
(d,0,\ldots,0)\in\ker\Phi,\qquad d>0.
\]

这里 $d$ 取为经过这些变换后能够出现的最小正首坐标。把任意关系的首坐标 $k$ 写成 $k=qd+r$，再减去 $q(d,0,\ldots,0)$；若 $r\ne0$，就会得到更小的正首坐标，矛盾。因此 $d\mid k$，特别地 $b_1$ 的阶正好是 $d$。

令 $B=\langle b_2,\ldots,b_t\rangle$。若 $k b_1\in B$，上面的整除性说明 $d\mid k$；于是 $\langle b_1\rangle\cap B=\{0\}$，而 $G$ 由这两个子群生成。由于 $G$ 阿贝尔，乘法映射 $\langle b_1\rangle\times B\to G$ 是同构。再对 $B$ 的生成元个数归纳，就得到有限个循环群的直积；$d=1$ 时删去平凡因子。

这段是课堂给出的存在性证明轮廓；初等因子和不变因子分解的唯一性仍采用教材中的有限生成 $\mathbb Z$-模结构定理。

\subsection{记忆点}


\begin{itemize}
\item $\mathbb Z$ 是无限但有限生成；有限生成不等于有限。
\item 循环群由一个元素决定，分类来自 $\mathbb Z$ 的核是 $0$ 还是 $m\mathbb Z$。
\item 外部直积把群放在有序对上；内部直积把同样的结构放回一个大群中。
\item 有限生成阿贝尔群的自由部分是 $\mathbb Z^r$，挠部分是有限循环群的直积。
\item 课堂证明用“最小关系 + 归纳”说明分解的存在性；唯一性通常用有限生成 $\mathbb Z$-模的结构定理。
\end{itemize}


\subsection{参考对照}


\begin{itemize}
\item [DF] §§5.1–5.2：Direct Products 与 Fundamental Theorem of Finitely Generated Abelian Groups。
\item [DS] §9.6：直积、内部直积与有限阿贝尔群结构定理。
\end{itemize}

\section{商群的三个同构定理与中心化子}


\textbf{本节用到的记号，先一次交代清楚}：

\begin{itemize}
\item $G$：群；$N\trianglelefteq G$、$H\trianglelefteq G$：正规子群；$A\le G$、$K\le G$：子群。
\item $G/N$：商群，元素是陪集 $gN$；$\cong$：同构；$[G:H]$：指数。
\item $HN=\{hn:h\in H,n\in N\}$：两个子群的乘积（$N$ 正规时它仍是子群）。
\item $H\cap N$：交集；$H/(H\cap N)\cong HN/N$（第二同构定理）。
\item $(G/N)/(H/N)\cong G/H$（第三同构定理，要求 $N\le H$）。
\item $C_G(A)=\{g\in G:ga=ag\ \text{对一切}\ a\in A\}$：$A$ 在 $G$ 中的中心化子；$C_G(x)$ 是 $A=\{x\}$ 的情形。
\item $Z(G)=C_G(G)$：群的中心。
\item $HK=G$ 且 $H\cap K=\{e\}$：内部直积的判据；此时 $G\cong H\times K$。
\end{itemize}


\subsection{中心与中心化子}


群 $G$ 的中心定义为

\[
Z(G)=\{z\in G:zg=gz\text{ 对所有 }g\in G\}.
\]

中心是 $G$ 的正规子群：若 $z\in Z(G)$，则对任意 $g\in G$ 有 $gzg^{-1}=z$。当 $G$ 阿贝尔时，$Z(G)=G$。

给定子集 $A\subseteq G$，其中心化子定义为

\[
C_G(A)=\{g\in G:ga=ag\text{ 对所有 }a\in A\}.
\]

特别地，$C_G(\{a\})$ 常简写为 $C_G(a)$。中心化子是子群；当 $A=G$ 时，$C_G(G)=Z(G)$。要注意：任意中心化子通常只是子群，不一定正规。

\subsection{第二同构定理}


设 $N\trianglelefteq G$，$H\le G$。则 $HN$ 是 $G$ 的子群，且

\[
H\cap N\trianglelefteq H,
\qquad HN/N\le G/N.
\]

第二同构定理给出

\[
H/(H\cap N)\cong HN/N.
\]

证明的自然映射是

\[
\theta:H\longrightarrow HN/N,
\qquad h\longmapsto hN.
\]

它是满同态，且

\[
\ker\theta=H\cap N.
\]

由第一同构定理即得结论。这个证明说明：第二同构定理本质上是把商映射 $G\to G/N$ 限制到 $H$ 上。

\subsection{第三同构定理}


设 $N\trianglelefteq G$、$H\trianglelefteq G$ 且 $N\le H$。则 $H/N\trianglelefteq G/N$，并且

\[
(G/N)/(H/N)\cong G/H.
\]

自然映射为

\[
\psi:G/N\longrightarrow G/H,
\qquad gN\longmapsto gH.
\]

它良定义，因为若 $gN=g'N$，则 $g^{-1}g'\in N\le H$，从而 $gH=g'H$。其核是 $H/N$，因此由第一同构定理得到上述同构。

记号容易混淆：如果把较小的正规子群记作 $N$，把较大的记作 $H$，公式就是 $(G/N)/(H/N)\cong G/H$；教材也可能交换字母，但包含关系必须保持。

\subsection{对应定理（第四同构定理）}


设 $N\trianglelefteq G$，令 $\pi:G\to G/N$ 为自然商映射。则

\[
\{K\le G:N\le K\}
\longleftrightarrow
\{\overline K\le G/N\}
\]

之间存在一一对应，具体为

\[
K\longmapsto K/N=\pi(K),
\qquad
\overline K\longmapsto \pi^{-1}(\overline K).
\]

这两个操作互为逆。对应保留以下结构：

\begin{itemize}
\item $K\trianglelefteq G$ 当且仅当 $K/N\trianglelefteq G/N$；
\item 若 $K\le L$，则 $[L:K]=[L/N:K/N]$；
\item 若另外有 $K\trianglelefteq G$，则中间商满足 $(G/N)/(K/N)\cong G/K$。
\end{itemize}


因此，商群把核 $N$ 压成单位元后，中间子群的格结构仍然完整保留下来。

\subsection{一张总图}


\[
\begin{array}{ccc}
K & \hookrightarrow & G\\[4pt]
\downarrow & & \downarrow\\[4pt]
K/N & \hookrightarrow & G/N
\end{array}
\]

这里 $N\le K\le G$；横向箭头表示包含，竖向箭头表示取商。反向操作是从 $G/N$ 的子群取原像，因而得到对应定理的一一对应。

\subsection{记忆点}


\begin{itemize}
\item 第二同构定理：$H$ 与 $N$ 的交控制 $H$ 在商群中的像。
\item 第三同构定理：连续商可以合并，先除 $N$ 再除 $H/N$ 等于直接除 $H$。
\item 对应定理：$G/N$ 的子群正好来自 $G$ 中包含 $N$ 的子群。
\item 中心 $Z(G)$ 是与整个群交换的元素；中心化子 $C_G(A)$ 是与指定集合 $A$ 交换的元素。
\end{itemize}


\subsection{参考对照}


\begin{itemize}
\item [DF] §3.3：第二、第三与第四（格）同构定理；中心化子见 §2.2。
\item [DS] §§9.5.15–9.5.16 对照第二、第三同构定理；中心化子按群论定义对照。
\end{itemize}

\section{同态、自同态、自同构与内自同构}


\textbf{本节用到的记号，先一次交代清楚}：

\begin{itemize}
\item $\operatorname{Hom}(G,H)$：$G$ 到 $H$ 的全体同态；$\operatorname{End}(G)=\operatorname{Hom}(G,G)$：自同态（复合下是幺半群）。
\item $\operatorname{Aut}(G)$：$G$ 的全体自同构（可逆自同态），在复合下成群；$\operatorname{Aut}(\mathbb Z_n)\cong\mathbb Z_n^\times$。
\item $c_g:G\to G,\ c_g(x)=gxg^{-1}$：由 $g$ 决定的共轭（内自同构）；$c_g\circ c_h=c_{gh}$。
\item $\operatorname{Inn}(G)=\{c_g:g\in G\}\trianglelefteq\operatorname{Aut}(G)$：内自同构群。
\item $\kappa:G\to\operatorname{Aut}(G),\ g\mapsto c_g$：内自同构同态；$\ker\kappa=Z(G)$，故 $\operatorname{Inn}(G)\cong G/Z(G)$。
\item $\operatorname{Out}(G)=\operatorname{Aut}(G)/\operatorname{Inn}(G)$：外自同构群。
\item $Z(G)$：中心；$G/Z(G)$：商群；$\cong$：同构；$\ker$：核。
\item $\varphi\cdot x=\varphi(x)$：自同构群作用于 $G$ 的写法（本节末尾的预告）。
\item $\mathbb Z_n^\times$：$\mathbb Z_n$ 的单位群，即与 $n$ 互素的剩余类。
\end{itemize}


\subsection{四个常用集合}


给定群 $G,H$，记

\[
\operatorname{Hom}(G,H)=\{\varphi:G\to H:\varphi\text{ 是群同态}\}.
\]

当定义域和值域相同时，得到

\[
\operatorname{End}(G)=\operatorname{Hom}(G,G),
\]

称为 $G$ 的自同态集合。可逆的自同态组成自同构群

\[
\operatorname{Aut}(G)=\{\varphi:G\to G:\varphi\text{ 是同构}\},
\]

运算是函数复合，单位元是恒等映射，逆元是逆函数。

在复合运算下，$\operatorname{End}(G)$ 通常只是含单位元的幺半群，$\operatorname{Hom}(G,H)$ 一般也没有自然的群结构；只有在额外的交换群或环结构下，才可以讨论点态加法。

\subsection{共轭与内自同构}


对每个 $g\in G$，定义共轭映射

\[
 c_g:G\to G,\qquad c_g(x)=gxg^{-1}.
\]

它是自同构，因为

\[
 c_g(xy)=gxyg^{-1}=(gxg^{-1})(gyg^{-1}).
\]

并且 $c_g^{-1}=c_{g^{-1}}$，以及

\[
c_g\circ c_h=c_{gh}.
\]

所有共轭映射组成的集合称为内自同构群：

\[
\operatorname{Inn}(G)=\{c_g:g\in G\}\le \operatorname{Aut}(G).
\]

映射

\[
\kappa:G\longrightarrow\operatorname{Aut}(G),
\qquad g\longmapsto c_g
\]

是群同态，其核为

\[
\ker\kappa=Z(G).
\]

因此由第一同构定理

\[
\operatorname{Inn}(G)\cong G/Z(G).
\]

在 $\operatorname{Aut}(G)$ 中，$\operatorname{Inn}(G)$ 是正规子群；商群 $\operatorname{Aut}(G)/\operatorname{Inn}(G)$ 常记为 $\operatorname{Out}(G)$，但本讲重点是内自同构本身。

正规性可由任意 $\alpha\in\operatorname{Aut}(G)$ 直接看出：

\[
\alpha c_g\alpha^{-1}=c_{\alpha(g)}\in\operatorname{Inn}(G).
\]

\subsection{例：$\operatorname{Aut}(\mathbb Z_n)$}


循环群 $\mathbb Z_n$ 由 $\bar 1$ 生成，因此一个同态 $\varphi:\mathbb Z_n\to\mathbb Z_n$ 完全由 $\varphi(\bar 1)=\bar k$ 决定，并且

\[
\varphi(\bar m)=\overline{km}.
\]

该映射是自同构，当且仅当 $\bar k$ 也是生成元；这等价于

\[
\gcd(k,n)=1.
\]

于是

\[
\operatorname{Aut}(\mathbb Z_n)\cong \mathbb Z_n^\times,
\]

其中 $\mathbb Z_n^\times$ 是模 $n$ 的可逆剩余类群，运算为乘法。这里“可逆”不是额外记忆点，而是保证乘法映射有逆映射的必要充分条件。

\subsection{与群作用的预告}


自同构群天然作用在 $G$ 上：$\varphi\cdot x=\varphi(x)$。共轭作用就是这个观点的一个特殊情况。下一节把“群元素对集合元素做变换”抽象成一般的群作用。

\subsection{记忆点}


\begin{itemize}
\item $\operatorname{End}(G)$ 允许不可逆的自映射；$\operatorname{Aut}(G)$ 只保留可逆者。
\item $g\mapsto c_g$ 的核是中心，所以 $G/Z(G)$ 描述了真正产生的共轭变换。
\item 循环群上的自同构由生成元的像决定；生成元恰好对应模 $n$ 的单位。
\end{itemize}


\subsection{参考对照}


\begin{itemize}
\item [DF] 关于 automorphism、inner automorphism 和 $\operatorname{Aut}(\mathbb Z_n)$ 的章节。
\item [DS] 群同态与群作用前的自同构例子。
\end{itemize}

\section{群作用}


\textbf{本节用到的记号，先一次交代清楚}：

\begin{itemize}
\item $G$：群；$X$（或 $\Omega$）：被作用的集合，不要求是群。
\item $G\curvearrowright X$ 或 $\rho:G\times X\to X,\ (g,x)\mapsto g\cdot x$：$G$ 在 $X$ 上的作用，满足 $e\cdot x=x$ 与 $(gh)\cdot x=g\cdot(h\cdot x)$。
\item $\rho_g:X\to X,\ \rho_g(x)=g\cdot x$：由 $g$ 给出的置换；$\rho:G\to\operatorname{Sym}(X)$ 是群同态。
\item $\operatorname{Sym}(X)$：$X$ 上全体双射构成的对称群。
\item $\ker\rho=\{g\in G:g\cdot x=x\ \text{对所有}\ x\}$：作用的核；核为 $\{e\}$ 时称作用\textbf{忠实}。
\item $G\cdot x=\{g\cdot x:g\in G\}$：$x$ 的轨道；$G_x=\{g:g\cdot x=x\}$：$x$ 的稳定子（下一节展开）。
\item $Z(G)$：中心；$\cong$：同构；$\le$：子群；$\times$：直积。
\item $GL_n(F)$ 与 $A\cdot v=Av$：矩阵群作用在列向量上的标准例子。
\end{itemize}


\subsection{定义}


设 $G$ 是群，$X$ 是集合。一个左群作用是一个映射

\[
\rho:G\times X\longrightarrow X,
\qquad (g,x)\longmapsto g\cdot x,
\]

满足

\[
 e\cdot x=x,
\qquad
(gh)\cdot x=g\cdot(h\cdot x)
\]

对所有 $g,h\in G$、$x\in X$ 成立。第二条说明群乘法与“先后作用”一致。

对每个 $g\in G$，定义置换

\[
\rho_g:X\to X,
\qquad \rho_g(x)=g\cdot x.
\]

由于 $g^{-1}$ 给出逆映射，$\rho_g\in\operatorname{Sym}(X)$。于是每个作用等价于一个群同态

\[
\rho:G\longrightarrow\operatorname{Sym}(X),
\qquad g\longmapsto\rho_g.
\]

作用的核是

\[
\ker\rho=\{g\in G:g\cdot x=x\text{ 对所有 }x\in X\}.
\]

若核为平凡群，则称作用是\textbf{忠实的}。因此“群作用”可以理解为把抽象群表示成一群具体的置换。

\subsection{基本例子}


\subsubsection{对称群的自然作用}


$\operatorname{Sym}(X)$ 作用在 $X$ 上：对 $\sigma\in\operatorname{Sym}(X)$ 和 $x\in X$，令 $\sigma\cdot x=\sigma(x)$。

\subsubsection{左正则作用}


群 $G$ 通过左乘作用在自身上：

\[
 g\cdot x=gx.
\]

该作用是忠实的，因为若 $g$ 固定所有 $x$，特别固定单位元 $e$，就有 $g=ge=e$。因此得到 Cayley 嵌入

\[
G\hookrightarrow\operatorname{Sym}(G).
\]

这说明每个群都同构于某个对称群的子群。

\subsubsection{共轭作用}


$G$ 作用在自身上，定义为

\[
 g\cdot x=gxg^{-1}.
\]

这个作用对应的同态就是 $g\mapsto c_g$，其核为 $Z(G)$。

\subsubsection{线性作用}


若 $G\le GL_n(F)$，则 $G$ 作用在向量空间 $F^n$ 上：

\[
 A\cdot v=Av.
\]

这是“群作用”把代数结构与几何对象连接起来的典型例子。

\subsection{作用的图像}


对固定 $x\in X$，作用给出一个轨道映射

\[
\alpha_x:G\longrightarrow X,
\qquad g\longmapsto g\cdot x.
\]

它的像是 $x$ 能够到达的所有位置；它的纤维将由稳定子陪集描述。这个观察是轨道--稳定子定理的核心。

\subsection{记忆点}


\begin{itemize}
\item 群作用的两条公理分别对应单位元和乘法相容。
\item 作用就是一个到对称群的同态；忠实性就是核平凡。
\item 左正则作用给出 Cayley 定理；共轭作用给出 $G/Z(G)\cong\operatorname{Inn}(G)$。
\end{itemize}


\subsection{参考对照}


\begin{itemize}
\item [DF] §1.7、§2.2：Group Actions、Stabilizers and Kernels。
\item [DS] §14.1.1–14.1.4：群作用、作用同态与 Cayley 定理。
\end{itemize}

\section{轨道、稳定子与轨道--稳定子定理}


\textbf{本节用到的记号，先一次交代清楚}：

\begin{itemize}
\item $G\curvearrowright X$：群作用；$g\cdot x$：$g$ 作用在 $x$ 上的结果。
\item $G\cdot x=\{g\cdot x:g\in G\}\subseteq X$：$x$ 的轨道；不同轨道要么相等要么不交，全体轨道构成 $X$ 的划分。
\item $G_x=\{g\in G:g\cdot x=x\}\le G$：$x$ 的稳定子。
\item $G/G_x$：稳定子的左陪集\textbf{集合}（$G_x$ 未必正规，所以这里一般不是商群）。
\item 轨道--稳定子定理：$G/G_x\to G\cdot x,\ gG_x\mapsto g\cdot x$ 是双射；$G$ 有限时 $|G\cdot x|=[G:G_x]=|G|/|G_x|$。
\item $[G:G_x]$：指数；$|X|=\sum_i|G\cdot x_i|$：按轨道求和。
\item $\operatorname{Sym}(X)$：$X$ 的对称群；$\cong$：同构。
\item $D_8$：正方形的对称群，本节用它演示轨道与稳定子。
\end{itemize}


\subsection{轨道与稳定子}


设 $G$ 作用在集合 $X$ 上，取 $x\in X$。$x$ 的轨道定义为

\[
G\cdot x=\{g\cdot x:g\in G\}\subseteq X.
\]

它表示从 $x$ 出发，在群作用下能够到达的全部位置。

$x$ 的稳定子（或固定子群）定义为

\[
G_x=\operatorname{Stab}_G(x)=\{g\in G:g\cdot x=x\}.
\]

这是 $G$ 的子群：单位元固定 $x$；若 $g,h$ 都固定 $x$，则 $gh$ 固定 $x$；若 $g$ 固定 $x$，则 $g^{-1}$ 也固定 $x$。

\subsection{轨道是等价类}


定义关系

\[
 x\sim y\quad\Longleftrightarrow\quad y=g\cdot x\text{ 对某个 }g\in G.
\]

这是等价关系：单位元给出自反性，取逆元给出对称性，作用的结合律给出传递性。因此不同轨道要么相等，要么不交，所有轨道构成 $X$ 的一个划分。

若只有一个轨道，即对任意 $x,y\in X$ 都存在 $g\in G$ 使 $g\cdot x=y$，则称作用是传递的。

\subsection{轨道与陪集}


考虑轨道映射

\[
\alpha_x:G\longrightarrow G\cdot x,
\qquad g\longmapsto g\cdot x.
\]

若 $g,h\in G$，则

\[
\alpha_x(g)=\alpha_x(h)
\Longleftrightarrow
g\cdot x=h\cdot x
\Longleftrightarrow
h^{-1}g\in G_x
\Longleftrightarrow
g\in hG_x.
\]

所以 $\alpha_x$ 的纤维正是稳定子 $G_x$ 的左陪集。于是它在陪集空间上诱导双射

\[
G/G_x\xrightarrow{\ \cong\ }G\cdot x,
\qquad gG_x\longmapsto g\cdot x.
\]

这里的 $G/G_x$ 只是左陪集集合；$G_x$ 一般不正规，所以不能把这个记号默认理解成商群。

\subsection{轨道--稳定子定理}


\textbf{定理。} 若 $G$ 有限，则

\[
|G\cdot x|=[G:G_x]=\frac{|G|}{|G_x|}.
\]

证明就是上面的陪集双射加拉格朗日定理。更一般地，即使 $G$ 无限，也有集合意义上的双射 $G/G_x\cong G\cdot x$；只有在有限情形下才能写成阶的商。

若作用传递且 $G,X$ 有限，则对任意 $x\in X$ 有

\[
|X|=|G\cdot x|=\frac{|G|}{|G_x|}.
\]

更一般地，若 $X$ 有限并取每条轨道的代表元 $x_1,\ldots,x_r$，则

\[
|X|=\sum_{i=1}^r|G\cdot x_i|
=\sum_{i=1}^r[G:G_{x_i}].
\]

\subsection{例：正方形的对称群}


令 $D_8$ 表示正方形的二面体群（按本笔记约定，阶为 $8$），它作用在四个顶点集合 $X$ 上。该作用是传递的。固定一个顶点 $x$，能保持 $x$ 不动的元素只有恒等变换和关于经过该顶点与中心的对角线反射，因此

\[
|G_x|=2,
\qquad |G\cdot x|=\frac{8}{2}=4=|X|.
\]

这个例子把“稳定子记录一个点还剩多少对称性”和“轨道记录这个点能移动到多少位置”同时展示出来。

\subsection{记忆点}


\begin{itemize}
\item 轨道是点在作用下的可达集合；稳定子是固定该点的群元素集合。
\item 轨道是等价类，所以轨道给出集合划分。
\item 轨道--稳定子定理的真正内容是 $G/G_x\cong G\cdot x$；有限公式只是其阶数版本。
\item 计算轨道大小时，先找稳定子，再用指数或拉格朗日定理。
\end{itemize}


\subsection{参考对照}


\begin{itemize}
\item [DF] §1.7、§2.2：orbits、stabilizers 与 coset bijection。
\item [DS] §14.1.5–14.1.9：轨道、稳定子和轨道--稳定子定理。
\end{itemize}


\section*{参考文献}
\begin{itemize}
\item \textbf{[DF]} D. S. Dummit, R. M. Foote. \emph{Abstract Algebra}, 3rd ed. §§1.7, 2.2, 4.4, 5.1--5.2. John Wiley \& Sons, 2004.
\item \textbf{[DS]} 邓小铁、王畅。《离散数学与结构》。未定稿，2023。§§9.5--9.6, 14.1。
\end{itemize}
\end{document}
